\section{Bases of normal tensors}

\begin{definition}[Basis of Normal Tensors (BNT)]\label{def:bnt}
    \lean{MPSTensor.IsBNT}
    \leanok
    \uses{def:normal, def:mpv}
    A \emph{basis of normal tensors} (BNT) for a total tensor $A_{\mathrm{tot}}$
    consists of $g$ block tensors $(A_1, \ldots, A_g)$ with bond dimensions
    $(D_1, \ldots, D_g)$ such that each $A_j$ is normal
    (Definition~\ref{def:normal}); at each system size~$N\geq1$ the MPV of
    $A_{\mathrm{tot}}$ lies in the span of the block MPVs
    $\{\ket{V^{(N)}(A_j)}\}_{j=1}^g$; and for sufficiently large~$N$ those block
    MPVs are linearly independent. Here normality means eventual block
    injectivity, as in Definition~\ref{def:normal}. Thus this is the algebraic
    form of \cite[Definition~IV.2]{Cirac2021Matrix}, rather than its primitive-channel
    formulation. In the canonical-form characterization of
    \cite[Section~2.3]{Cirac2016MPDO_arXiv}, a BNT is a minimal family of
    representatives for $\widetilde A_k^i=e^{i\phi_k}X_kA_j^iX_k^{-1}$ among
    the normal blocks.
\end{definition}

\begin{definition}[Basis of normal tensors, coefficient form]\label{def:bnt_cpsv}
    \lean{MPSTensor.IsCPSVBasisOfNormalTensors}
    \leanok
    \uses{def:nt_cpsv, def:mpv}
    % Source: Cirac2017MPDO_AnnPhys \S~II.C, lines 271--274.
    The coefficient form of a basis of normal tensors: each block $A_j$ is a
    CPSV normal tensor (Definition~\ref{def:nt_cpsv}), the MPV of $A$ is given at
    every positive length by explicit coefficients $c_j^{(N)} \in \C$ with
    $\ket{V^{(N)}(A)} = \sum_{j=1}^{g} c_j^{(N)} \ket{V^{(N)}(A_j)}$, and the
    block MPVs $\{\ket{V^{(N)}(A_j)}\}$ are eventually linearly independent.
    This is the literal coefficient-form definition of
    \cite[Section~2.3]{Cirac2016MPDO_arXiv}. It differs from
    Definition~\ref{def:bnt} only in the chosen normality predicate.
    Theorem~\ref{thm:cpsv_bnt_is_bnt} proves the implication between the two
    forms.
\end{definition}

\begin{theorem}[Distinct BNT members are not gauge-phase equivalent]
    \label{thm:cpsv_bnt_blocks_not_gaugePhaseEquiv}
    \lean{MPSTensor.IsCPSVBasisOfNormalTensors.blocks_not_gaugePhaseEquiv}
    \leanok
    \uses{def:bnt_cpsv, def:blocks_not_gauge_phase_equiv}
    Distinct members of a basis of normal tensors are not gauge-phase
    equivalent after identifying equal bond dimensions.
\end{theorem}

\begin{proof}\leanok
    At every sufficiently large length the matrix product vectors of the basis
    members are linearly independent. A gauge-phase equivalence between two
    distinct members would give
    $\ket{V^{(N)}(A_k)}=\zeta^N\ket{V^{(N)}(A_j)}$ for a nonzero scalar
    $\zeta$. Thus one vector would be a scalar multiple of the other,
    contradicting linear independence.
\end{proof}

% CPSV16 Proposition 2.7 (`prop:char-BNT`), lines 278--280 and its
% Appendix A restatement at lines 1137--1142.
% Provenance: source-label audit completed in issue #5917.
% Gap: docs/paper-gaps/cpsv16_bnt_characterization_active_blocks.tex.
\begin{theorem}[Printed BNT characterization]
    \label{thm:cpsv_prop27_printed_status}
    \notready
    \uses{def:bnt_cpsv, def:cpsv_canonical_form, def:gauge_phase_equiv, thm:cpsv_bnt_characterization}
    Let
    \begin{align}
        A^i & =\bigoplus_{k=1}^{r}\mu_k\widetilde A_k^i
        \notag
    \end{align}
    be a tensor in canonical form.  According to
    \cite[Proposition~2.7]{Cirac2016MPDO_arXiv}, a family $(A_j)_j$ is a basis
    of normal tensors for $A$ if and only if every listed normal block
    $\widetilde A_k$ is gauge-phase equivalent to some $A_j$ and no two
    distinct $A_j$ are gauge-phase equivalent.

    This is false when a coefficient $\mu_k$ vanishes.  Such a block is
    invisible in every positive-length matrix product vector and need not occur
    in a basis recovered from those vectors.  If only the active indices $k$
    with $\mu_k\ne0$ are required to be represented, the characterization
    holds by Theorem~\ref{thm:cpsv_bnt_characterization}.
\end{theorem}

\begin{theorem}[Characterization by canonical-form representatives]
    \label{thm:cpsv_bnt_characterization}
    \lean{MPSTensor.isCPSVBasisOfNormalTensors_iff_canonicalForm_covered_and_minimal}
    \lean{MPSTensor.SectorDecomposition.%
        exists_phase_match_of_isCPSVBasisOfNormalTensors}
    \leanok
    \uses{def:bnt_cpsv, def:nt_cpsv, def:block_diagonal_tensor, def:same_mpv2_pos, def:blocks_not_gauge_phase_equiv}
    Let $A$ have a canonical-form MPV expansion
    \begin{align}
        \ket{V^{(N)}(A)}
         & = \sum_{k=1}^{r}\mu_k^N\ket{V^{(N)}(\widetilde A_k)}
        \notag
    \end{align}
    at every positive length, where the tensors $\widetilde A_k$ are normal
    tensors of positive bond dimension and every coefficient $\mu_k$ is
    nonzero. Let $(A_j)_{j=1}^{g}$ be tensors of
    positive bond dimension. Then
    $(A_j)_{j=1}^{g}$ is a basis of normal tensors for $A$ if and only if:
    \begin{enumerate}
        \item every $A_j$ is normal;
        \item for every $k$ there are an index $j$, an
              invertible matrix $X_k$, and a phase $e^{i\phi_k}$ such that
              $\widetilde A_k^i=e^{i\phi_k}X_kA_j^iX_k^{-1}$ for every
              physical index $i$;
        \item no two distinct tensors $A_j$ and $A_{j'}$ are related in this
              way.
    \end{enumerate}
    This is \cite[Proposition~2.7]{Cirac2016MPDO_arXiv}, read with the
    nonzero-coefficient convention of Definition~\ref{def:cpsv_canonical_form}.
\end{theorem}

\begin{proof}\leanok
    \uses{thm:cpsv_normal_mpv_phase_gauge, thm:cpsv_normal_mpv_phase_alternative, thm:cpsv_bnt_blocks_not_gaugePhaseEquiv, lem:bnt_finite_index_overlap_orthonormal_li, lem:finite_geometric_sum_eventual_vanishing, lem:sector_decomposition_coeff_not_eventually_zero, def:mpv_phase_class_data, thm:cpsv_normal_separated_eventually_li, thm:mpv_decomposition}
    Group the summands into the MPV
    phase classes of Definition~\ref{def:mpv_phase_class_data}, and choose one
    representative from each class. By
    Theorem~\ref{thm:cpsv_bnt_blocks_not_gaugePhaseEquiv}, eventual linear
    independence of a BNT excludes gauge-phase equivalence between two
    distinct $A_j$.

    Suppose that some phase-class representative is not gauge-phase equivalent
    to any $A_j$. Partition all phase classes into the unmatched classes $T$
    and the remaining classes. Every remaining representative is a nonzero
    scalar multiple of some $\ket{V^{(N)}(A_j)}$ at each positive length.
    Distinct unmatched representatives have vanishing mutual overlaps, and
    their overlaps with every $A_j$ also vanish. Explicitly, for $t\ne t'$ in
    $T$ and every $j$,
    \begin{align}
        O_{\widetilde A_t\widetilde A_{t'}}(N) & \longrightarrow0,
        \notag                                                     \\
        O_{\widetilde A_t A_j}(N)              & \longrightarrow0.
        \notag
    \end{align}
    By Lemma~\ref{lem:bnt_finite_index_overlap_orthonormal_li}, the combined
    family
    \begin{align}
        \{\ket{V^{(N)}(\widetilde A_t)}:t\in T\}
        \cup \{\ket{V^{(N)}(A_j)}:1\leq j\leq g\}
        \notag
    \end{align}
    is linearly independent for all sufficiently large $N$. Comparing the
    canonical-form expansion with the BNT expansion gives a relation
    \begin{align}
        \sum_{t\in T}c_{N,t}\ket{V^{(N)}(\widetilde A_t)}
        +\sum_{j=1}^{g}\alpha_{N,j}\ket{V^{(N)}(A_j)}
         & =0.
        \notag
    \end{align}
    Linear independence forces every $c_{N,t}$ to vanish. Hence the
    coefficient of every unmatched phase class vanishes for all sufficiently
    large $N$.
    For an unmatched class $t$, let $q$ range over its canonical-form
    summands, write $\mu_q$ for the corresponding canonical-form weight, and
    write $\zeta_q$ for the phase relating that summand to the chosen
    representative. Then
    $c_{N,t}=\sum_{q\in t}(\zeta_q\mu_q)^N$.
    Every base $\zeta_q\mu_q$ is nonzero because every coefficient is nonzero.
    Eventual vanishing of this finite geometric sum would imply vanishing at
    $N=0$, where its value is the positive number of summands in the class.
    This contradiction proves coverage. The normal-tensor phase theorem turns
    the resulting MPV phase relation into the stated gauge-phase relation.

    Conversely, coverage gives, for each canonical-form summand,
    $\ket{V^{(N)}(\widetilde A_k)}
        =\zeta_k^N\ket{V^{(N)}(A_{j(k)})}$. Hence
    \begin{align}
        \ket{V^{(N)}(A)}
         & = \sum_j
        \left(\sum_{k:j(k)=j}(\mu_k\zeta_k)^N\right)
        \ket{V^{(N)}(A_j)}.
        \notag
    \end{align}
    This is the required BNT span at every positive length. Pairwise
    gauge-phase distinction and normality give eventual linear independence,
    completing the converse.
\end{proof}

